Formalization is in progress. The companion uses Mathlib's actual upper half-plane and \texttt{Complex.UnitDisc}; the latter is the open unit ball in $\mathbb C$. The shared Cayley maps use the fractions in the manuscript, with nonvanishing denominators and the inequalities placing their values in the correct domains proved explicitly. The resulting equivalence proves bijectivity between these domains; it is not yet registered as a metric-space isometry or a smooth equivalence.

For \hyperref[hyperbolic.cayley]{the Cayley transform}, Lean proves the formula for $1-|C(z)|^2$, the imaginary part of $C^{-1}(\zeta)$, both inverse identities, and the actual complex derivatives using \texttt{HasDerivAt}. For the separate \hyperref[hyperbolic.disk-metric]{disk-metric proposition}, the proof of the inverse metric coefficients substitutes that derivative into the previously defined half-plane tangent energy and proves the value $4|v|^2/(1-|\zeta|^2)^2$. This follows the manuscript's cancellation of the $|1-\zeta|$ factors. Identifying this with the metric on a smooth tangent bundle, proving the path-length and distance equivalence, and deriving the actual area measure remain further steps. The separate \hyperref[hyperbolic.disk-rotations]{rotation proposition} and \hyperref[hyperbolic.disk-geodesics]{disk-geodesic classification} have not yet been formalized.

The shared definition of the \hyperref[hyperbolic.point-pair]{point-pair invariant} is the exact quotient $|z-w|^2/(4\operatorname{Im}z\operatorname{Im}w)$, using complex norm squared. Lean proves the \hyperref[hyperbolic.point-pair-properties]{elementary properties}: nonnegativity, symmetry, vanishing exactly on equal points, and invariance under the actual $G$-action. Invariance reuses the previously checked two-point ratio $J$ after proving $u=J/4$. The inverse Cayley difference identity and the \hyperref[hyperbolic.point-pair-disk]{disk formula for $u$} follow the prose's subtraction and denominator cancellation, with disk membership supplying the needed positivity.

The specialization $u(i,w)=|C(w)|^2/(1-|C(w)|^2)$ is now also checked. The Lean proof substitutes the disk points $C(i)$ and $C(w)$ into the general disk formula, uses both inverse Cayley identities, and simplifies $C(i)=0$, just as in the expanded orbit proof. This proves the substitution step; constructing a rotation and the resulting element $g_2^{-1}kg_1$ remains a separate obligation.

The \hyperref[hyperbolic.point-pair-distance]{distance corollary}, relating $u$ to the distance defined as an infimum of path lengths, remains open. The \hyperref[hyperbolic.pair-orbits]{ordered-pair orbit classification}---including the explicit element $g_2^{-1}kg_1$ and the zero-radius case---and the separate \hyperref[hyperbolic.radial-functions]{invariant-function corollary} are also not yet Lean theorems. These claims are not replaced by definitions that assume the desired conclusions.

For the separate \hyperref[hyperbolic.polar-coordinates]{polar-coordinate proposition}, \hyperref[hyperbolic.ball-area]{disk-area corollary}, and \hyperref[hyperbolic.radial-integration]{radial integration proposition}, remaining targets are the radius relation $\rho=\tanh(r/2)$ with $r=d_{\mathbb H}(i,w)$, the polar metric and genuine area measure, the ball-area formula, and the radial integral with factor $4\pi$. That integral must prove both the nonnegative case (including $+\infty$) and absolute integrability for the complex-valued case. The elementary coordinate formulas already implemented do not establish those integration claims.

Machine evidence for the implemented portions is recorded in \texttt{lean/audits/chapter01-release-check-2026-11-12/record.json}, against its source hashes. The section and chapter remain in progress; Steven's correspondence review is pending independently of compilation and the axiom audit.
